% !TeX root = ./chapter_2_main.tex
 
% \tocomplete{...} marks technical details that still need to be filled in.
% Remove or replace each occurrence once the corresponding information is available.
\providecommand{\tocomplete}[1]{\textbf{[TODO: #1]}}
 
\section{Injection into the black fiber}
\label{sec:black_fiber_injection}
 
The second optical field required for the conveyor-belt configuration is delivered by a conventional polarization-maintaining single-mode fiber, referred to throughout this work as the "black fiber". It counterpropagates with respect to the $1064,mathrm{nm}$ field guided through the Hollow-Core Photonic Crystal Fiber (HCPCF). Their interference generates the moving optical lattice described in Sec.~\ref{sec:conveyor_belt}. Efficient coupling, a well-defined spatial mode, and stable polarization are therefore required to obtain a deep and reproducible conveyor-belt potential.

The black-fiber branch was constructed as an extension of the Optical Dipole Trap (ODT) system described already existing. The injection strategy was first validated on a dedicated test bench and subsequently transferred to the definitive optical branch. 

The black fiber is operated at $1064\,\mathrm{nm}$ and has a nominal mode-field radius of approximately $5.2\,\mu\mathrm{m}$. Its input tip includes a glass endcap approximately $300\,\mu\mathrm{m}$ long, which improves power handling but modifies the effective optical distance between the focusing element and the guided mode. Specifically, the fiber is a polarization-maintaining, single-mode cable manufactured by Sch\"{a}fter+Kirchhoff (model PMC-780-1-18-xxx / PMC-780-5.3-NA012-1-APC-xxx-P, housed in a $900\,\mu\mathrm{m}$ buffer). It features an effective numerical aperture ($\mathrm{NA}_{e^2}$) of $0.08 \pm 0.005$ at $780\,\mathrm{nm}$, low-stress FC/APC ($8^\circ$) connectors on both ends with the polarization slow axis aligned to the connector index key, and a specified operating wavelength range from the $770\,\mathrm{nm}$ cut-off up to a maximum of $1100\,\mathrm{nm}$.

\subsection{Gaussian-beam characterization procedure}
\label{subsec:gaussian_beam}

Throughout this work, the size and divergence of the optical beam are
characterized using a common beam-profiling procedure. The same method is
applied before and after the fiber whenever the injection geometry, input
beam parameters, or output wavefront have to be determined.
The beam under test is first aligned horizontally at a fixed height above
the optical table. This allows its transverse intensity distribution to be
sampled at a sequence of longitudinal positions without introducing
additional steering optics between measurements. A camera mounted on a
translation stage is then moved along the propagation direction, and an
image of the beam cross-section is recorded at each position.

Each recorded image is cropped to isolate the beam spot from background
light and stray reflections. The resulting intensity distribution is fitted
with a two-dimensional Gaussian function, from which the beam radii along the
two transverse axes are extracted. Repeating this analysis at several
positions provided the beam radius as a function of the longitudinal
coordinate $z$.
For each transverse direction, the measured beam radius is fitted to the
standard Gaussian-beam propagation law
\begin{equation}
    w(z) =
    w_0
    \sqrt{
        1+
        \left(
            \frac{z-z_0}{z_R}
        \right)^2
    },
\end{equation}
where $w_0$ is the waist radius, $z_0$ is the waist position, and $z_R$ is
the associated Rayleigh range. This fit provides the size and position of the
beam waist and allowed the wavefront curvature at any relevant plane to be
determined. The reconstructed beam parameters could then be compared with
the target fiber-mode parameters or used to select and position the
injection optics.

\subsection{Injection setup}
\label{subsec:black_fiber_setup}

The injection setup for the black fiber comprises the pre-existing ODT injection setup, shared with the HCPCF branch, and a branch built specifically for the black fiber downstream of PBS2.
The HCPCF branch was previously built, but a HWP was added before the focusing lens for the polarization control required to implement the conveyor belt.

The pre-existing setup, shown in Figure~\ref{fig:HC_pdf}, consists of:
\begin{itemize}
    \item HWP1 and the optical isolator, protecting the laser from back-reflections;
    \item HWP2 and PBS1, regulating the optical power delivered to the rest of the setup;
    \item HWP3 and PBS2, splitting the beam between the HCPCF branch and the black-fiber branch;
    \item on the HCPCF branch, fed by the first output of PBS2: an AOM operating at $110\,\mathrm{MHz}$, used as a fast switch, with the first diffraction order used downstream and the zeroth order blocked by a beam dump
    \item a $\sim3\times$ beam expander, collimating and enlarging the beam
    \item a mirror (M1) and a dichroic mirror, both on tilt stages, used for the fiber-injection alignment, with the dichroic mirror also combining the $1064\,\mathrm{nm}$ beam with a $780\,\mathrm{nm}$ beam used for atom characterization inside the HCPCF core
    \item a focusing lens with $f=50\,\mathrm{mm}$, on a horizontal micrometric translation stage
    \item a HWP, for the polarization adjustment before the HCPCF, mounted on a rotation stage
    \item the HCPCF on an $x$-$y$ translation stage terminated with an FC/APC connector
\end{itemize}

\begin{figure}[h!]
    \centering
    \includegraphics[width=0.8\textwidth]{Cap_2_figures/HC.pdf}
    \caption{Optical scheme of the pre-existing ODT injection setup, from HWP1 to the HCPCF, with the second output of PBS2 leading to the black-fiber branch described below. \tocomplete{verify AOM orientation in the schematic}}
    \label{fig:HC_pdf}
\end{figure}

Downstream of PBS2, the branch built for the black fiber, shown in Figure~\ref{fig:black_pdf}, consists of:
\begin{itemize}
    \item an acousto-optic modulator (AOM), providing fast control of the beam intensity and frequency (\tocomplete{manufacturer, model, diffraction order used, drive frequency});
    \item a steering mirror, directing the beam into the telescope;
    \item a telescope formed by an $f=-30\,\mathrm{mm}$ lens and an $f=150\,\mathrm{mm}$ lens (\tocomplete{lens separation}), re-collimating the beam to a radius of approximately $1.6\,\mathrm{mm}$, the target value derived in Section~\ref{subsec:mode_matching_endcap};
    \item two steering mirrors, directing the beam toward the fiber-injection assembly;
    \item a half-wave plate, setting the incident linear polarization along the slow axis of the fiber (\tocomplete{HWP model and alignment procedure}, Figure~\ref{fig:polarization_maintaining});
    \item an $f=25\,\mathrm{mm}$ achromatic doublet (Thorlabs AC127-025-C), mounted on a micrometric translation stage along the propagation axis, focusing the beam into the fiber. The axial position of the lens is used to compensate for the displacement of the fiber-mode plane caused by the fiber endcap, discussed in Section~\ref{subsec:mode_matching_endcap};
    \item the black fiber, mounted on the branch through its output connector.
\end{itemize}

\tocomplete{manufacturer and model; standard telecom-fiber designation; total fiber length; nominal mode-field diameter and numerical aperture at $1064\,\mathrm{nm}$; orientation of the slow axis relative to the connector key; input and output connector types and angle; endcap geometry; spectral operating range; distance from the AOM to the telescope; lens mount; translation-stage model, resolution and travel; nominal lens-to-fiber distance; fiber connector and mount; diagnostic pick-offs or power-measurement points along the branch}

\begin{figure}[h!]
    \centering
    \includegraphics[width=0.9\textwidth]{Cap_2_figures/black_fiber.pdf}
    \caption{Optical scheme of the black-fiber branch downstream of PBS2.}
    \label{fig:black_pdf}
\end{figure}

Figure~\ref{fig:black_setup_photo} shows a photograph of the complete injection setup.
\begin{figure}[h!]
    \centering
    \includegraphics[width=0.8\textwidth]{Cap_2_figures/black_setup_photo.jpeg}
    \caption{Photograph of the complete injection branch. The black-fiber branch is on the left, while the pre-existing ODT injection setup is on the right.}
    \label{fig:black_setup_photo}
\end{figure}

\begin{figure}[h!]
    \centering
    \includegraphics[width=0.4\textwidth]{Cap_2_figures/polarization_maintaining.png}
    \caption{Orientation of the input linear polarization relative to the slow and fast axes of the polarization-maintaining fiber.}
    \label{fig:polarization_maintaining}
\end{figure}


\subsection{Beam-profile characterization method}
\label{subsec:beam_char_black_fiber}

The size and divergence of the beams involved in the black-fiber injection were characterized with the same beam-profiling approach used in Section~\ref{subsec:gaussian_beam}.

For each transverse direction, the measured radius was fitted to the Gaussian-beam propagation law, this fit gave the size and position of the beam waist, which could then be compared with the target fiber-mode parameters or used to select and position the injection optics. The procedure was repeated both upstream and downstream of the fiber, after each modification or addition of the injection optics, to obtain the beam parameters used in the mode-matching calculations of Section~\ref{subsec:mode_matching_endcap}.

\subsection{Mode matching and the fiber endcap}
\label{subsec:mode_matching_endcap}

Coupling a Gaussian beam into a single-mode fiber requires matching the focused field to the fundamental fiber mode in waist size, waist position and transverse alignment. The standard approach uses a collimated input beam and a fixed collimating lens that focuses it directly onto the fiber facet, and this was the strategy first used for the black fiber, then substituted by a f25 shifting lens. The mode-matching parameters used in the definitive branch (Section~\ref{subsec:black_fiber_setup}) were derived and validated on a dedicated test bench (Fig. \ref{fig:test_black}), independent of the complete ODT apparatus and not including the AOM or the telescope. The test bench consisted of:
\begin{itemize}
    \item HWP1 and the optical isolator, as in Section~\ref{subsec:black_fiber_setup};
    \item HWP2 and PBS1 (with a beam dump on the rejected port), regulating the optical power delivered to the rest of the setup;
    \item HWP3 and PBS2, splitting the beam between the HCPCF and black-fiber branches;
    \item Mirror 1, directing the beam reflected by PBS2 towards the beam expander;
    \item an adjustable beam expander (ZBE3C2X-8X Achromatic Zoom Beam Expander, AR Coated: 1050 nm-1650 nm), setting the diameter and residual divergence of the beam sent to the black fiber;
    \item HWP4 and PBS3 (with a beam dump on the rejected port), for additional power control and polarization cleaning;
    \item Mirror 2 and Mirror 3, routing the beam towards the fiber injection stage;
    \item HWP5, adjusting the polarization axis before injection into the polarization-maintaining fiber;
    \item a focusing element coupling the beam into the black fiber: initially a fixed $f=18\,\mathrm{mm}$ fiber collimator (Sch\"{a}fter+Kirchhoff, 60FC-4-A18-03), later replaced by the $f=25\,\mathrm{mm}$ achromatic doublet (Thorlabs AC127-025-C) on a micrometric translation stage (z-stage) also used in the definitive branch;
\end{itemize}

Injecting a nominally collimated beam directly through the fixed $f=18\,\mathrm{mm}$ collimator gave an injection efficiency of approximately $20\%$. Introducing additional divergence upstream, equivalent to inserting a weak $f=-1000\,\mathrm{mm}$ CHI LHA DETTO lens before the collimator, raised the efficiency above $50\%$. This indicates that the plane of the fiber mode does not coincide with the nominal focal plane of the collimator.

The beam profile after the fiber was measured (Figure~\ref{fig:black_fiber_characterization}a) and its reconstructed waist compared with the mode-field diameter (MFD) specified by the manufacturer at $1064\,\mathrm{nm}$ (Figure~\ref{fig:black_fiber_characterization}b). The measured waist agreed with the specified MFD, confirming that the transmitted light corresponds to the fundamental fiber mode.

\begin{figure}[htbp]
    \centering
    % Sottofigura (a)
    \begin{subfigure}[b]{0.49\textwidth}
        \centering
        \includegraphics[width=\textwidth]{Cap_2_figures/beam_profile_after.png}
        \caption{Beam-radius measurements after the black fiber and fit to the Gaussian-beam propagation law of Equation~\eqref{eq:gaussian_propagation_black_fiber}. \tocomplete{measurement axis, fit parameters and uncertainties, optical configuration}}
        \label{fig:beam_profile_after}
    \end{subfigure}
    \hfill
    % Sottofigura (b)
    \begin{subfigure}[b]{0.49\textwidth}
        \centering
        \includegraphics[width=\textwidth]{Cap_2_figures/manufacturer_MFD.png}
        \caption{Manufacturer-specified mode-field diameter of the black fiber at $1064\,\mathrm{nm}$, used as the target value in the mode-matching calculation below.}
        \label{fig:manufacturer_MFD}
    \end{subfigure}
    
    \caption{Black fiber beam characterization and manufacturer specifications.}
    \label{fig:black_fiber_characterization}
\end{figure}

The beam profile analysis shown in Fig.\ref{fig:black_fiber_characterization}a uses the standard Gaussian beam propagation model. The beam radius evolution $w(z)$ along the propagation axis $z$ for both the $x$ and $y$ transverse directions is described by:

\begin{align}
    w_x(z) &= w_{0x}\sqrt{1 + \left(\frac{z - z_{0x}}{z_{Rx}}\right)^2} \, , \quad \text{with} \quad z_{Rx} = \frac{\pi w_{0x}^2}{M_x^2 \lambda} \label{eq:fit_wx} \\
    w_y(z) &= w_{0y}\sqrt{1 + \left(\frac{z - z_{0y}}{z_{Ry}}\right)^2} \, , \quad \text{with} \quad z_{Ry} = \frac{\pi w_{0y}^2}{M_y^2 \lambda} \label{eq:fit_wy}
\end{align}

Based on the fit of the experimental data, the extracted parameters for the $x$-axis are:
\begin{align*}
    z_{0x} &= 3.97 \pm 0.43\,\mathrm{mm} \\
    w_{0x} &= (5.67 \pm 0.10) \times 10^{-3}\,\mathrm{mm} \\
    M_x^2 &= 1.02 \, (\text{fixed})
\end{align*}

For the $y$-axis, the extracted fit parameters are:
\begin{align*}
    z_{0y} &= 3.85 \pm 0.48\,\mathrm{mm} \\
    w_{0y} &= (5.42 \pm 0.11) \times 10^{-3}\,\mathrm{mm} \\
    M_y^2 &= 1.02 \, (\text{fixed})
\end{align*}

The goodness of the fit is confirmed by the high coefficient of determination for both axes:
\begin{align*}
    R_x^2 &= 0.9993 \\
    R_y^2 &= 0.9991
\end{align*}

The mismatch between the nominal collimator focal plane and the fiber-mode plane is attributed to the endcap spliced onto the fiber input, a glass segment approximately $300\,\mu\mathrm{m}$ long used to improve power handling (Figure~\ref{fig:fiber_endcap}). The endcap modifies the effective optical path between the external focusing lens and the fiber facet, displacing the focused waist relative to the position expected for a bare fiber tip. Standard fiber collimators are designed assuming the fiber mode lies at their nominal focal plane, so this displacement is not compensated automatically and must be corrected in the injection design. Reproducing the correction with a fixed, long-focal-length diverging lens would be impractical to position and would add aberrations to the injected beam.
A focusing lens mounted on a micrometric translation stage, giving direct control of the axial position of the focused waist, was used instead.

\begin{figure}[h!]
    \centering
    \includegraphics[width=0.4\textwidth]{Cap_2_figures/fiber_endcap.png}
    \caption{Geometry of the endcap spliced onto the black-fiber input. \tocomplete{connector orientation, propagation direction}}
    \label{fig:fiber_endcap}
\end{figure}

For a collimated beam of radius $w_{\mathrm{in}}$ focused by a lens of focal length $f$, Equation~\eqref{eq:gaussian_waist_lens} gives the resulting waist radius $w_{\mathrm{f}}$. Solving for the input radius required to obtain a target $w_{\mathrm{f}}$,
\begin{equation}
    w_{\mathrm{in}} = \frac{\lambda f}{\pi w_{\mathrm{f}}}.
    \label{eq:required_input_radius}
\end{equation}
Using $\lambda=1.064\,\mu\mathrm{m}$, $f=25\,\mathrm{mm}$, and $w_{\mathrm{f}}=5.2\,\mu\mathrm{m}$, taken as half of the manufacturer-specified MFD (Figure~\ref{fig:manufacturer_MFD}), gives
\begin{equation}
    w_{\mathrm{in}}
    =
    \frac{
        (1.064\times10^{-6}\,\mathrm{m})
        (25\times10^{-3}\,\mathrm{m})
    }{
        \pi
        (5.2\times10^{-6}\,\mathrm{m})
    }
    \approx
    1.6\,\mathrm{mm}.
\end{equation}
The beam expander was set to produce this radius at the injection lens, and the resulting beam was profiled before injection to verify both its size and its residual divergence.

The axial position of the $f=25\,\mathrm{mm}$ lens was calibrated by propagating light in the reverse direction: with $7\,\mathrm{W}$ injected through the fixed $f=18\,\mathrm{mm}$ collimator and collected at the output by the $f=25\,\mathrm{mm}$ lens, the lens position was adjusted until the output beam was collimated. This confirmed that the lens had been placed at the axial position corresponding to the effective fiber-mode plane, and fixed the lens position later reproduced in the definitive branch.

\begin{figure}[h!]
    \centering
    \includegraphics[width=0.8\textwidth]{Cap_2_figures/test_black.pdf}
    \caption{Setup of the black-fiber injection test bench, used to validate the mode-matching procedure and calibrate the axial position of the $f=25\,\mathrm{mm}$ lens. At the fiber output the polarization stability was characterized with the setup of Section~\ref{sec:black_fiber_polarization}.}
    \label{fig:test_black}
\end{figure}

On the test bench, with the $f=25\,\mathrm{mm}$ lens at the calibrated position and forward propagation, the power measured after the black fiber and the $f=18\,\mathrm{mm}$ output collimator reached approximately $70\%$ of the power measured before the fiber, directly after the $f=25\,\mathrm{mm}$ lens. 

With the definitive branch of Section~\ref{subsec:black_fiber_setup}, after adding the AOM and the telescope, the measured injection efficiency reached approximately $75\%$,confirming that the mode-matching procedure transfers to the complete optical layout.

\section{Polarization stability characterization}
\label{sec:black_fiber_polarization}

After optimizing the injection efficiency, the polarization stability of the light transmitted through the black fiber was characterized. This quantity is relevant to the optical conveyor belt because polarization fluctuations are converted into optical-power fluctuations by the polarization-sensitive elements of the final setup, affecting the stability of the resulting optical potential. The two orthogonal linear-polarization components at the fiber output were measured simultaneously; their time-domain stability was analyzed with the Allan deviation, and their frequency content with the power spectral density, obtained from the Fourier transform of the recorded signals and compared with the expected shot-noise level.

\subsection{Detection setup and polarization alignment}

The polarization-analysis setup is shown in Figure~\ref{fig:test_black}. The beam emerging from the black fiber was collimated and sent through a half-wave plate (HWP) followed by a polarizing beam splitter (PBS), whose transmitted and reflected components were detected simultaneously by two photodiodes.

During the initial alignment, the transmitted component was measured with a home-made photodiode connected to channel 1, and the reflected component with a Thorlabs PDA100A2 connected to channel 2, operated at $0\,\mathrm{dB}$ internal gain with $6\,\mathrm{dB}$ of additional electrical attenuation. The total optical power incident on the PBS was $4.2\,\mathrm{mW}$.

The HWP was adjusted to obtain approximately equal power at the two PBS outputs, corresponding to an input polarization close to $45^\circ$ relative to the PBS axes. At this operating point, a small polarization rotation transfers power between the two outputs, producing anticorrelated variations in the two detector signals, which gives high sensitivity to polarization fluctuations.

Before the stability measurements, the PBS orientation was optimized to ensure that its finite extinction ratio did not limit the measurement. Rotating the incident linear polarization and finely adjusting the PBS orientation gave polarization extinction ratios of $31\,\mathrm{dB}$ in reflection and $33\,\mathrm{dB}$ in transmission.


\tocomplete{output-collimator focal length; HWP and PBS models; acquisition-device model; photodiode connections and termination impedances; acquisition bandwidth and voltage ranges}

\subsection{Alignment to the fiber principal axes}

The polarization analyzer was referenced to the principal axes of the polarization-maintaining fiber, using an input HWP before the fiber and, after the fiber, an output collimator, a steering mirror, a second HWP, the PBS, and a power meter.

The two HWPs were adjusted to maximize the power transmitted through the PBS while minimizing the reflected power and its fluctuations. This condition corresponds to exciting one principal axis of the fiber and aligning the emerging linear polarization with one PBS axis; the optimal input-HWP angle was $131^\circ$.

The output HWP was adjusted to obtain a maximum transmitted power of $2.68\,\mathrm{mW}$. With the power meter moved to the reflection port, the reflected signal was minimized by jointly tuning the PBS orientation and the two HWPs, giving a minimum reflected power of $4\,\mu\mathrm{W}$, including a background of $1.5\,\mu\mathrm{W}$, against a maximum reflected power of approximately $2\,\mathrm{mW}$. Without background subtraction, the corresponding polarization extinction ratio is
\begin{equation}
    \mathrm{PER}
    =
    10\log_{10}
    \left(
        \frac{P_{\mathrm{max}}}{P_{\mathrm{min}}}
    \right)
    =
    10\log_{10}
    \left(
        \frac{2\,\mathrm{mW}}{4\,\mu\mathrm{W}}
    \right)
    \approx 27\,\mathrm{dB}.
\end{equation}

As a further check, the fiber transmission efficiency was measured again, with input and output powers of $4.18\,\mathrm{mW}$ and $2.88\,\mathrm{mW}$, respectively:
\begin{equation}
    \eta_{\mathrm{fiber}}
    =
    \frac{2.88\,\mathrm{mW}}{4.18\,\mathrm{mW}}
    \approx 69\%,
\end{equation}
consistent with the efficiency measured during injection optimization.

\subsection{Photodiode calibration}

The two detector outputs were converted from voltage to equivalent optical power before comparing the polarization components.

The home-made detector is based on a BPW34 photodiode followed by an OP27G transimpedance amplifier with a gain of $10\,\mathrm{k\Omega}$. Using a photodiode responsivity at $1064\,\mathrm{nm}$ of
\begin{equation}
    \mathcal{R}_{\mathrm{HM}}
    \approx
    0.7\,\mathrm{A/W}\times0.3
    =
    0.21\,\mathrm{A/W},
\end{equation}
the resulting voltage responsivity is
\begin{equation}
    G_{\mathrm{HM}}
    =
    \mathcal{R}_{\mathrm{HM}}G_{\mathrm{TIA}}
    \approx
    2.1\,\mathrm{V/mW}.
\end{equation}
The detector output stage has a $1\,\mathrm{k\Omega}$ output impedance; the corresponding loading correction was included where required by the acquisition-input impedance.

For the Thorlabs PDA100A2, a photodiode responsivity of $0.4\,\mathrm{A/W}$ at $1064\,\mathrm{nm}$ was used, together with the selected internal gain and any additional electrical attenuation, to convert the recorded voltage into optical power.

\tocomplete{measured offsets and dark signals; acquisition-input impedance; PDA100A2 transimpedance gain at $0\,\mathrm{dB}$ and $10\,\mathrm{dB}$; uncertainty of both voltage-to-power conversion factors}

\subsection{Data acquisition and normalized polarization signal}

For the stability measurements, the home-made detector was placed on the reflected port and connected to channel 1, and the Thorlabs PDA100A2 on the transmitted port and connected to channel 2, operated at $10\,\mathrm{dB}$ internal gain without external attenuation.

For each measurement, $10^5$ samples were acquired simultaneously from both channels. Traces of $600\,\mathrm{ms}$, $60\,\mathrm{s}$, and $600\,\mathrm{s}$ total duration were recorded to probe different timescales; the $60\,\mathrm{s}$ trace was acquired both in standard sampling mode, preserving faster fluctuations, and in High-Resolution mode, improving the resolution of slower fluctuations.

After background subtraction and detector calibration, the measured voltages were converted into the transmitted and reflected powers $P_{\mathrm{T}}(t)$ and $P_{\mathrm{R}}(t)$. The total detected power is
\begin{equation}
    P_{\Sigma}(t)
    =
    P_{\mathrm{T}}(t)+P_{\mathrm{R}}(t),
\end{equation}
and the normalized differential signal is
\begin{equation}
    D(t)
    =
    \frac{
        P_{\mathrm{T}}(t)-P_{\mathrm{R}}(t)
    }{
        P_{\mathrm{T}}(t)+P_{\mathrm{R}}(t)
    }.
    \label{eq:normalized_polarization_difference}
\end{equation}
\tocomplete{verify that Equation~\eqref{eq:normalized_polarization_difference} matches the normalization used in the analysis code, e.g. $D(t)=P_{\mathrm{T}}(t)/\overline{P_{\mathrm{T}}}-P_{\mathrm{R}}(t)/\overline{P_{\mathrm{R}}}$}
A fluctuation of the total transmitted power appears mainly in $P_{\Sigma}(t)$, while a polarization fluctuation transfers power between the two PBS outputs and appears predominantly in $D(t)$; the normalization suppresses common-mode power fluctuations to first order.

\tocomplete{sampling frequency for each trace; Nyquist frequency and frequency resolution; acquisition-device (oscilloscope or DAQ) model; anti-aliasing filter; exact definition of High-Resolution mode}

\subsection{Allan-deviation analysis}

The time-domain stability of the polarization signal was quantified with the overlapping Allan deviation. For a uniformly sampled sequence $D_k$ with sampling interval $\tau_0$, averages were computed over adjacent groups of $m$ samples,
\begin{equation}
    \overline{D}_j(\tau)
    =
    \frac{1}{m}
    \sum_{k=jm}^{(j+1)m-1}D_k,
    \qquad
    \tau=m\tau_0,
\end{equation}
and the Allan variance from the difference between consecutive averages,
\begin{equation}
    \sigma_D^2(\tau)
    =
    \frac{1}{2}
    \left\langle
        \left[
            \overline{D}_{j+1}(\tau)
            -
            \overline{D}_{j}(\tau)
        \right]^2
    \right\rangle,
    \qquad
    \sigma_D(\tau) = \sqrt{\sigma_D^2(\tau)}.
\end{equation}

The Allan deviation gives the typical change in the average polarization signal between consecutive intervals of duration $\tau$, and identifies the timescales over which averaging improves the polarization stability and those over which slow drift dominates. A decreasing $\sigma_D(\tau)$ indicates that averaging suppresses the dominant fluctuations; a minimum identifies the optimum averaging time; an increase at longer $\tau$ indicates the onset of drift or other low-frequency instabilities. The slope of $\sigma_D(\tau)$ on a logarithmic scale distinguishes different classes of noise.

The Allan deviation was evaluated for the individual polarization channels, for the total-power signal $P_{\Sigma}(t)$, and for the normalized differential signal $D(t)$, to determine whether the observed instability is dominated by common-mode optical-power noise or by a redistribution of power between the two polarization components.

\tocomplete{overlapping or non-overlapping estimator used; removal of mean or linear drift; range of averaging times; confidence intervals or number of averages at each $\tau$; treatment of the trace end points}

\subsection{Frequency-domain analysis and comparison with the shot-noise limit}

The frequency content of the measured power fluctuations was analyzed with the power spectral density (PSD) of the calibrated photodiode signal, computed from its discrete Fourier transform after subtracting the mean value of each trace. The PSD describes how the variance of the measured optical power is distributed over frequency.

\tocomplete{window function; direct FFT or Welch method; segment length and overlap; sampling frequency and frequency resolution; single- or double-sided PSD convention; whether a linear trend was removed; units used in the final plots ($\mathrm{W^2/Hz}$ or $\mathrm{W/\sqrt{Hz}}$)}

The measured spectrum was compared with the shot-noise limit expected for a coherent optical field at the same average power. For a photodiode with responsivity $\mathcal{R}$ receiving an average optical power $\overline{P}$, the mean photocurrent is $\overline{I}=\mathcal{R}\,\overline{P}$, and the single-sided shot-noise PSD of the photocurrent is
\begin{equation}
    S_{I,\mathrm{SN}}
    =
    2e\overline{I}
    =
    2e\mathcal{R}\,\overline{P},
\end{equation}
where $e$ is the elementary charge. The equivalent optical-power PSD follows by dividing by the squared responsivity,
\begin{equation}
    S_{P,\mathrm{SN}}
    =
    \frac{S_{I,\mathrm{SN}}}{\mathcal{R}^2}
    =
    \frac{2e\overline{P}}{\mathcal{R}}.
    \label{eq:shot_noise_power_psd}
\end{equation}
Using $\mathcal{R}=\mathrm{QE}\,e/(\hbar\omega)$, this can be written as
\begin{equation}
    S_{P,\mathrm{SN}}
    =
    \frac{2\hbar\omega\,\overline{P}}{\mathrm{QE}},
\end{equation}
where $\mathrm{QE}$ is the photodiode quantum efficiency. For unit quantum efficiency, this reduces to the quantum limit
\begin{equation}
    S_{P,\mathrm{QL}}
    =
    2\hbar\omega\,\overline{P}.
    \label{eq:quantum_shot_noise_limit}
\end{equation}

Equation~\eqref{eq:quantum_shot_noise_limit} was used as the reference shot-noise level for comparison with the measured optical-power PSD, both expressed with the same single-sided spectral convention and in units of $\mathrm{W^2/Hz}$. A measured spectrum above this limit indicates excess technical noise in addition to the fundamental photon shot noise. If the spectra are instead plotted as amplitude spectral densities, the corresponding quantum limit is
\begin{equation}
    \sqrt{S_{P,\mathrm{QL}}}
    =
    \sqrt{2\hbar\omega\,\overline{P}},
\end{equation}
with units of $\mathrm{W/\sqrt{Hz}}$.